\section{Discussion}
\label{sec:discussion}

\paragraph{OxfordPet outlier.}
OxfordPet shows a flat/clean ratio of $204$, far above any other dataset (next highest: Cars196 at $13.6$).
This means the AS$^1$ perturbation is disproportionately large relative to the task gradient scale,
likely because OxfordPet has tiny per-class gradient magnitudes due to its small class count
(37 classes $\to$ 4 classes per task).
A dataset-adaptive $\lambdaflat$ or $\rho$ that scales with the task gradient norm would address this;
we leave adaptive perturbation as future work.

\paragraph{ImageNet-A: LR sensitivity, not method failure.}
The $\mathrm{lr}=0.01$ setting yields FAA~$=45.999$, $13.9$pp below EWC-LoRA.
With $\mathrm{init\_lr}=0.04$, $\mathrm{lr}=0.02$ (warm-start cosine), FAA improves to $58.92$
($-0.97$pp vs.\ EWC-LoRA CNN), with NME-head FAA~$=64.45$ surpassing the EWC-LoRA reference.
The root cause is that ImageNet-A has 10 tasks of 20 classes each with only 20 epochs:
a $\mathrm{lr}=0.01$ schedule underfits each task, amplifying the plasticity deficit
that the Fisher penalty exacerbates.
A moderate warm-start ($\mathrm{init\_lr}=0.04$) restores sufficient task-level learning
without destabilizing the Fisher constraint.

\paragraph{When does Cars196 not benefit?}
Cars196 shows a slightly negative $\Delta$FAA ($-0.46$pp) but positive $\Delta$AAA ($+0.85$pp).
The flat/clean ratio is $13.6$---larger than typical---while fisher/clean is $0.19$,
suggesting the Fisher constraint is active but the flatness perturbation may be pulling the
adapter toward regions that are flat but not optimal for the final task.
A lower $\lambdaflat$ or smaller $\rho$ for this dataset would likely recover the FAA improvement.

\paragraph{Relationship to EWC-LoRA.}
EWC-LoRA~\citep{zheng2025ewclora} is the stability component of FS-LoRA without flatness control.
The comparison shows that adding decoupled AS$^1$ to EWC-LoRA consistently improves
performance on datasets where the Fisher penalty is lightly active (Flowers, CUB200, Aircraft),
and leaves performance similar when Fisher dominates.
This is consistent with the orthogonality principle: AS$^1$ adds a genuinely new signal,
not a redundant one.

\paragraph{When does orthogonality break?}
Orthogonality rests on spectral separation between current-task Hessian and past-task Fisher.
This separation weakens when:
(1)~Tasks share similar class distributions or visual concepts, reducing the geometric difference
between $H_t$ and $F_{t'}$.
(2)~The LoRA rank is very large, making the adapter space approach the full parameter space
where the spectral gap shrinks.
(3)~The backbone is partially unfrozen, mixing backbone and adapter gradients.
Future work should characterize conditions under which $|\cosff|$ grows and whether
the method degrades gracefully.

\paragraph{Connection to theory.}
The orthogonality finding connects to the sequential PAC-Bayesian analysis
in~\citet{Anonymous2026PaperB}, which shows that both the KL complexity and
sharpness terms in the continual generalization bound reduce to the admissible update
subspace $\mathcal{W}_{\Delta,t}$.
FS-LoRA implements both terms within $\DeltaW$-space, consistent with the theory's
prescription that the relevant flatness and stability measures should be computed
in the adapter geometry.
The orthogonality finding is a \emph{new} empirical and mechanistic contribution
beyond what the theory establishes.

\paragraph{Dual ascent: principled but passive.}
Section~\ref{sec:dual_ascent} presents dual ascent as the outer-loop Lagrangian update
on $\lambdaewc$ (Eq.~\eqref{eq:dual_ascent}).
On all benchmarks in Section~\ref{sec:experiments},
$\Dtilde^{\mathrm{Fisher}}$ remained below $\delta$ throughout,
consistent with the orthogonality principle: a near-orthogonal flatness update
does not drive the adapter toward Fisher-sensitive directions.
The dual ascent mechanism is inactive in all reported results but provides
a principled adaptive extension for task sequences with high semantic overlap.
